\section{敏感性与情景讨论}
\subsection{返航储备参数}
问题一的返航储备敏感性已在图~\ref{fig:q1sens} 中给出。随着预留电量增加，可供载荷使用的能量减少，平均安全载荷呈非增趋势。该趋势与约束式~\eqref{eq:q1payload} 一致；但本文只报告参考结果所含 10\%--30\% 的离散情景，不对区间外进行外推。

\subsection{分区数变化}
问题四的 $K=2$ 与 $K=3$ 对比表明，分区数增加不必然降低总资源需求。独立组不能共享资源，且任务区划改变并发重叠结构；在给定固定排程与库存下，$K=3$ 增加一组 C 型电池缺口，使总缺口从 3 增至 4。此结论针对用户提供的新版 Q4 TEX 所述分区目标与峰值口径，不代表任意时序重排后的总体最优结论。

\subsection{模型边界情景}
若将天气、临时禁飞区或通信衰落引入模型，应相应扩展可行架次集合和链路可用状态；若允许各区之间临时借调资源，则需增加转场时间和跨组调度变量。以上属于后续扩展，本稿不赋予数值预测。
